%! Probability Models — Unit 3, Lesson 3.2 — Homework
% The lesson's SCORED (1 to 5) individual practice and the LAST component in the packet. Assigned
% at Close & Assign, done outside class. PACKET form: every final answer is a \gradedblank.
% THIRD CONTEXT: the notes teach on a food-truck festival, the Guided Practice runs on a sports
% roster, and these items are on a bookstore loyalty program.
% Items 3 and 4 are the deliberate CONTRAST PAIR: the whole circle against the "only" region.
% Item 5 is the crux head-on (adding the two circle totals). Item 1 interprets the expected count;
% item 5's written part is also the justification; item 6 is the spiral (Lesson 3.1 relative frequency).
% THE DATA
%   Wheel, 10 equal slots: $1 off 1-5; $3 off 6-8; bookmark 9; book 10.
%     P($3 off) = 3/10; 500 spins -> 150; P(not $3 off) = 7/10.
%   Members (60): mystery 26, graphic 17, both 7; only mystery 19, only graphic 10; union 36;
%     neither 24; the sum 43. P(mystery) = 26/60 = 13/30; P(only mystery) = 19/60.
%   Spiral (60): teen print 12, teen e-book 6, adult print 21, adult e-book 21; print 33,
%     e-book 27; teens 18, adults 42. Print: 33/60 = 0.55. Print among teens: 12/18 = 0.67.
\documentclass[12pt]{article}
\usepackage{saar-article}
\usepackage{saar-boxes}
\usepackage{grading}
\gradingassignment{saar-u03-l02-hw}
\gradingcourse{stats_analysis_algebraic_reasoning}
\newcommand{\TallMath}[1]{$\displaystyle #1\rule[-1.4em]{0pt}{3.2em}$}
% Handwriting room inside every \begin{work} block. Moves the blank and the
% key together, so the two can never drift apart.
\setlength{\workrowsep}{6pt}

\begin{document}

\pageheader{Unit 3, Lesson 3.2}{Homework\strut}
{\small\itshape Write each final answer in its box. Only what is inside a box is graded.}\par
% NAMESTRIP: no \namedateperiod here. The student writes their name once, on the
% cover this component is stapled behind.

\vspace{0.08in}

% Authored identically in homework/ and homework_key/.
\begin{remindbox}
\small
\textbf{This is your graded homework.} It is scored 1 to 5 and due at the start of the first
class after two study halls. It is assigned today and done on your own, outside class --- flip
back to the \emph{Keep in Mind} box on the cover and to rows 2 and 4 of your notes.
\end{remindbox}

\vspace{0.06in}

\boxguard[12]
\begin{scenariobox}[Turning Page Books: The Loyalty Program]{cerulean}
At checkout, a loyalty member spins a rewards wheel with $10$ equal slots: slots $1$--$5$ give
\$1 off, slots $6$--$8$ give \$3 off, slot $9$ gives a free bookmark, and slot $10$ gives a free
book.
\end{scenariobox}

\vspace{0.06in}

\boxguard
\begin{notesbox}{The Rewards Wheel}
Every answer is about the wheel. Write probabilities as fractions in simplest form.
\begin{enumerate}[leftmargin=1.6em, itemsep=6pt, topsep=3pt]

  \item Find $P(\text{\$3 off})$ for one spin. \hfill \blankbox{\gradedblank[id=q1a, type=short, points=1, answer={3/10}, check=expression]{2.2cm}{13.5pt}}

        \vspace{3pt}
        The store expects $500$ spins next month. How many spins should land on \$3 off, and
        what does that prediction mean?
        \gradedresponse{\gradedblank[id=q1b, type=rubric, points=1, display={About 150 spins, since (3/10)(500) = 150. It is a prediction, not a promise: the spins are random, so the real count should be near 150 but may not be exactly 150.}]{\linewidth}{2in}}

  \item Use the complement rule to find $P(\text{not \$3 off})$. \hfill \blankbox{\gradedblank[id=q2a, type=short, points=1, answer={7/10}, check=expression]{2.2cm}{13.5pt}}

        \vspace{3pt}
        Are ``\$1 off'' and ``\$3 off'' complementary, mutually exclusive only, or neither? \par\vspace{2pt}\hfill \blankbox{\gradedblank[id=q2b, type=short, points=1, answer={mutually exclusive}, check=text, display={mutually exclusive only}]{5.4cm}{13.5pt}}

\end{enumerate}
\end{notesbox}

\boxguard[30]
\begin{notesbox}{The Members Survey}
The store surveyed its $60$ members about last month. Pick one member at random. Probabilities: fractions in simplest form.

\begin{center}
\begin{tikzpicture}[scale=0.74]
  \draw[thick, charcoal] (-6.0,-1.85) rectangle (6.0,2.15);
  \begin{scope}[blend mode=multiply]
    \fill[frost] (-1.05,0) circle[radius=1.6];
    \fill[goldbg] (1.05,0) circle[radius=1.6];
  \end{scope}
  \draw[thick, cerulean] (-1.05,0) circle[radius=1.6];
  \draw[thick, goldacc] (1.05,0) circle[radius=1.6];
  \node[font=\footnotesize\bfseries, align=center] at (-4.05,1.4) {Mystery\\($26$ in all)};
  \node[font=\footnotesize\bfseries, align=center] at (4.05,1.4) {Graphic novel\\($17$ in all)};
  \node at (-1.75,0) {$19$};
  \node at (0,0) {$7$};
  \node at (1.75,0) {$10$};
  \node at (4.7,-1.3) {$24$};
\end{tikzpicture}
\end{center}
\vspace{-6pt}
\begin{enumerate}[leftmargin=1.6em, itemsep=5pt, topsep=2pt, start=3]

  \item How many bought a mystery? \hfill \blankbox{\gradedblank[id=q3a, type=short, points=1, answer={26}, check=numeric]{1.6cm}{13.5pt}}
        \quad $P(\text{mystery})$: \; \blankbox{\gradedblank[id=q3b, type=short, points=1, answer={13/30}, check=expression]{2.2cm}{13.5pt}}

  \item How many bought \textbf{only} a mystery? \hfill \blankbox{\gradedblank[id=q4a, type=short, points=1, answer={19}, check=numeric]{1.6cm}{13.5pt}}
        \quad $P(\text{only a mystery})$: \; \blankbox{\gradedblank[id=q4b, type=short, points=1, answer={19/60}, check=expression]{2.2cm}{13.5pt}}

  \item Priya adds the circle totals: \emph{``$26 + 17 = 43$ members bought a mystery or a
        graphic novel.''} How many really did? \hfill \blankbox{\gradedblank[id=q5a, type=short, points=1, answer={36}, check=numeric]{1.6cm}{13.5pt}}
        \par\vspace{2pt}
        Why is Priya's count too high? Name the members she counted twice.
        \gradedresponse{\gradedblank[id=q5b, type=rubric, points=1, display={The 7 members who bought both are in each circle total, so adding the totals counted them twice: 43 - 7 = 36. Counting each region once: 19 + 7 + 10 = 36.}]{\linewidth}{2in}}

  \item \emph{Back to Lesson 3.1.} The same $60$ members, by age and the format they prefer.
        Write each answer as a decimal to two places.

        \begin{center}
        \renewcommand{\arraystretch}{1.15}
        \begin{tabular}{@{} l c c c @{}}
        \toprule
        & \textbf{Print} & \textbf{E-book} & \textbf{Total} \\
        \midrule
        Teen & 12 & 6 & 18 \\
        Adult & 21 & 21 & 42 \\
        \textbf{Total} & 33 & 27 & 60 \\
        \bottomrule
        \end{tabular}
        \end{center}
        \vspace{-3pt}
        Relative frequency of members who prefer print: \hfill \blankbox{\gradedblank[id=q6a, type=short, points=1, answer={0.55}, check=numeric, tolerance=0.01]{1.8cm}{13.5pt}}
        \par\vspace{2pt}
        Among the teens, relative frequency who prefer print: \hfill \blankbox{\gradedblank[id=q6b, type=short, points=1, answer={0.67}, check=numeric, tolerance=0.01]{1.8cm}{13.5pt}}

\end{enumerate}
\end{notesbox}

\boxguard[5]
\begin{spiralbox}
\textbf{Next class: Lesson 3.3, Addition and Multiplication Rules.} Today you counted ``or''
region by region; next class that count becomes a rule that works with probabilities, and
multiplying along a tree's branches comes with it.
\end{spiralbox}

\end{document}
